\subsection*{Part A: Pull-back geodesics}

\noindent We train a VAE \cite{kingma2022autoencodingvariationalbayes} with 2D latent space on a 3-class MNIST subset (2048 observations, Gaussian prior and likelihood) and compute geodesics under the pull-back metric induced by the decoder mean $f(z)$. Geodesics minimise the discrete curve energy $\mathcal{E}(c)=\sum_i\|f(z_{i+1})-f(z_i)\|^2$ via L-BFGS, initialised with linear interpolation between fixed endpoints (see code snippets) \cite{hauberg2026latent}.

\autoref{fig:part_a} shows 25 geodesics for two independent training runs. The learned representations differ a bunch -- classes arrange differently due to the identifiability problem inherent in the VAE: the model can reparametrise the latent space without changing the data distribution, so different training runs recover geometrically distinct but equally valid embeddings \cite{hauberg2026latent}. In both runs, geodesics curve along the data manifold, going around low density gaps, and are on average ${\sim}2.5\times$ longer than the straight-line Euclidean distance. While geodesic lengths are theoretically invariant to such reparametrisations, we still observe variation across seeds (e.g.\ pair 1687--130: 17.6 vs.\ 19.1), indicating that different seeds converge to slightly different manifold approximations $f$. This is what directly motivates the ensemble approach in Part B.

\begin{figure}[H]
    \centering
    \begin{subfigure}{0.48\linewidth}
        \centering
        \includegraphics[width=\linewidth]{media/latent_space_seed_42.png}
        \caption{Training with seed 42}
        \label{fig:latent_seed_42}
    \end{subfigure}
    \hfill
    \begin{subfigure}{0.48\linewidth}
        \centering
        \includegraphics[width=\linewidth]{media/latent_space_seed_69.png}
        \caption{Training with seed 69}
        \label{fig:latent_seed_69}
    \end{subfigure}
    \caption{Latent space with 25 geodesics for two training runs. The class layout rotates across seeds, but geodesics consistently follow the data manifold in both cases.}
    \label{fig:part_a}
\end{figure}